\clearpage
\section{数据保存与操作衔接}
\label{sec:design-consistency}
接受纸条时，用户写好回信并确认后，服务端要改纸条状态、建会话并写入首轮消息。这些记录需要一起保存，避免用户重新打开页面时看到不完整的会话。表\ref{tab:design-transactions}列出了几种类似的操作。

\begin{table}[H]
\centering\small
\caption{关键业务操作的关联保存与恢复结果}\label{tab:design-transactions}
\begin{tabularx}{\linewidth}{@{}p{28mm}p{42mm}X@{}}
\toprule 业务操作 & 关联保存内容 & 操作衔接方式 \\ \midrule
注册与登录 & 账号与当前会话 & 注册形成个人空间，登录进入已有空间 \\
发现候选 & 本次对象与待决定卡片 & 继续查看同一张卡片，完成选择后再抽取 \\
收藏或移除 & 收藏、冷却与决定结果 & 收藏用于后续联系，移除后暂缓再次遇见 \\
接受真人纸条 & 纸条、会话与首轮消息 & 接受并回信后，转入双方持续聊天 \\
确认生成图 & 情绪图、完成状态与预览处理 & 保存选择的形象，继续用于聊天展示 \\
历史兼容处理 & 既有数据与必要转换结果 & 沿用已有记录，接续当前使用流程 \\
\bottomrule
\end{tabularx}
\end{table}

已经集中保存的操作会一起提交关联更新，保存失败则恢复原状态。图片生成依赖外部服务，所以先记录进度并提供预览，用户确认后再保存正式立绘。保存成功后再次进入设置页，页面会读取已保存的立绘。

中断后重新进入，可以继续看已有卡片、聊天历史和素材。遗留的图片生成任务会标为失败，是否重试由用户决定。已有记录通过兼容处理继续读取。重置互动会重新初始化虚拟体验和默认展示设置，账号、真人交流、匹配收藏和媒体仍保留。

后续工作包括完善关联更新的提交范围、异步结果确认和中断恢复，并补充备份及公开使用所需的管理功能。
